\section{Experiment 6: The Hidden Cost of Printing (Formatting \& Serial Latency)}

\subsection{Aim}
To measure the computational overhead of converting numerical representations into human-readable ASCII text in memory, determine the physical transmission latency of UART communication at 115200 baud, and contrast I/O delays against pure CPU arithmetic throughput.

\subsection{Background}
Converting binary values to text constitutes a representation transformation that requires substantial algorithmic processing.
\begin{itemize}[noitemsep]
    \item \textbf{Decimal Formatting:} Transforming an integer to decimal text requires repeated modulo and division operations by 10 (base-10 radix conversion), which cannot be implemented using simple binary bit shifts.
    \item \textbf{Floating-Point Formatting:} Formatting a floating-point number into ASCII requires separating the sign, biased exponent, and mantissa, scaling by powers of 10, performing rounding, and emitting character buffers (e.g., via the Dragon4 or Grisu algorithm).
    \item \textbf{Hexadecimal Formatting:} Converting to hexadecimal is inherently fast because each hex character maps directly to a 4-bit nibble, requiring only bit shifts and masking (\texttt{\& 0x0F}).
    \item \textbf{UART Framing at 115200 Baud:} Physical asynchronous serial communication over UART uses an 8-N-1 frame: 1 start bit, 8 data bits, no parity bit, and 1 stop bit, totaling \textbf{10 physical bit periods per character}.
\end{itemize}

\subsection{Prediction (Recorded Prior to Measurement)}
% TODO(Team): Record your hypothesis here BEFORE running the firmware.
% Expected snprintf cost per format specifier; expected UART transmission
% latency (theoretical calc is fine here since it's arithmetic, not a measurement).
\textit{[Pending: team hypothesis not yet recorded.]}

\subsection{Method}
We benchmarked in-memory conversions using \texttt{snprintf()} across 50 iterations over 10 trials to isolate pure CPU conversion time without UART transmission. We then transmitted a single ASCII character (\texttt{'A'}) over the hardware UART, calling \texttt{Serial.flush()} to ensure complete transmission out of the hardware FIFO, timing the duration using \texttt{micros()} and \texttt{ESP.getCycleCount()}.

\subsection{Observations}
% Data source: ../data/exp6_printing.txt (raw Serial Monitor capture from speed_of_numbers.ino, command '6')
\textit{[Pending: no hardware measurements recorded yet. Run Experiment 6 on the ESP32 and populate \texttt{data/exp6\_printing.txt}.]}

\subsection{Analysis}
\textit{[Pending: analysis withheld until measurements are collected.]}

\subsection{What We Learnt}
\textit{[Pending.]}

\subsection{LLM Log}
\begin{llmbox}[LLM Interaction Record: UART Serial Baud Rate Math]
\textbf{Prompt to LLM:} \textit{[Pending: record the verbatim prompt given to the LLM.]} \\
\textbf{LLM Response Summary:} \textit{[Pending.]} \\
\textbf{Board Agreement:} \textit{[Pending.]}
\end{llmbox}
